% !TeX program = xelatex
\documentclass[UTF8,fontset=mac,zihao=-4]{ctexart}
\usepackage[a4paper,margin=28mm]{geometry}
\usepackage{amsmath,amssymb,amsthm}
\usepackage[hidelinks]{hyperref}
\setlength{\parindent}{2em}
\setlength{\parskip}{0.35em}
\newcommand{\energy}{\mathcal{E}}
\newcommand{\variation}[2]{\delta #1[#2]}
\newtheorem{observation}{观察}
\title{从变分到方程\\[0.5em]\large 一份关于能量方法的写作札记}
\author{QuillTeX 1.1.0 示例文稿}
\date{}
\pagestyle{plain}
\begin{document}
\maketitle
\begin{abstract}
从一个带固定端点的能量泛函出发，沿着试探函数、驻值条件和弱形式，
把一维边值问题的几种描述联系起来。本文保留计算的中间步骤，
也记录写作时需要核对的假设、符号与交叉引用。
\end{abstract}
% !TeX root = ../main.tex
\section{先写下能量}
\label{sec:energy}
设 $f\in C([0,1])$。在满足 $u(0)=u(1)=0$ 的函数中，考虑
\begin{equation}
  \energy[u]=\int_0^1\left(\frac12 |u'(x)|^2-f(x)u(x)\right)\,\mathrm{d}x.
  \label{eq:energy}
\end{equation}
第一项衡量函数变化的代价，第二项描述外部作用。
我们先在光滑函数上计算，再说明表达式为何自然延伸到弱导数的框架。

\begin{observation}
固定端点意味着允许的扰动也必须在两个端点为零。
这是后续分部积分中边界项消失的原因。
\end{observation}

取满足 $v(0)=v(1)=0$ 的试探函数 $v$，沿 $u+tv$ 展开，得到
\begin{align*}
  \energy[u+tv]-\energy[u]
    &=t\int_0^1(u'v'-fv)\,\mathrm{d}x
      +\frac{t^2}{2}\int_0^1|v'|^2\,\mathrm{d}x.
\end{align*}
这个展开同时给出一阶变化和余项，而不必预先猜测微分方程。

% !TeX root = ../main.tex
\section{驻值条件怎样变成方程}
\label{sec:stationary}
记 $\variation{\energy}{u}(v)$ 为方向 $v$ 上的一阶变化。
若 $u$ 是极小点，则所有允许方向的一阶变化都为零，即
\begin{equation}
  \variation{\energy}{u}(v)
    =\int_0^1(u'v'-fv)\,\mathrm{d}x=0.
  \label{eq:weak}
\end{equation}
当 $u$ 足够光滑时，分部积分把式~\eqref{eq:weak} 写成
$\int_0^1(-u''-f)v\,\mathrm{d}x=0$。于是得到边值问题
\begin{equation}
  -u''=f\quad\text{于 }(0,1),\qquad u(0)=u(1)=0.
  \label{eq:boundary}
\end{equation}
反过来，满足式~\eqref{eq:weak} 的函数具有
\[
  \energy[u+v]-\energy[u]=\frac12\int_0^1|v'|^2\,\mathrm{d}x\geq 0.
\]
因此驻值条件在这里已经足以保证极小性；固定端点还排除了非零常数扰动。

\begin{observation}
式~\eqref{eq:weak} 只含一阶导数。把试探函数空间取为 $H_0^1(0,1)$，
就得到同一问题的弱形式，而无需要求未知函数具有连续的二阶导数。
\end{observation}

% !TeX root = ../main.tex
\section{用一个例子核对计算}
\label{sec:review}
令 $f\equiv 1$，式~\eqref{eq:boundary} 的解为
\[
  u(x)=\frac{x(1-x)}{2}.
\]
它在两个端点取零，且 $-u''=1$。
代入式~\eqref{eq:energy} 可得 $\energy[u]=-1/24$，
与驻值计算一致。

整理文稿时，需要分别核对三个层次：允许函数的边界条件、
一阶变化中的符号，以及弱形式与经典方程之间所需的正则性。
第~\ref{sec:energy} 节和第~\ref{sec:stationary} 节的引用，
把这些检查重新连接到推导本身。

\end{document}
